\begin{abstract}
%Cyber threat intelligence (CTI) analysts routinely convert noisy, unstructured security artifacts into standardized, automation-ready outputs. This mapping is still brittle when performed by general-purpose language models. Prior CTI LLM work has largely relied on supervised fine-tuning (SFT). At the same time, CTI standards and community-maintained resources provide structured targets and curated mappings that enable deterministic, programmatic verification and thus verifiable rewards. We explore reinforcement learning with verifiable rewards (RLVR) for CTI and introduce Minerva, a unified dataset and training pipeline spanning diverse CTI subtasks with verifiers that score canonical identifier and structured output predictions. We also address reward sparsity, where many prompts receive zero reward under rollout. We propose a lightweight self-training mechanism that generates additional verified trajectories and distills them back into the model using the original prompts. Experiments across several LLM backbones show consistent gains in accuracy and robustness across multiple CTI benchmarks and tasks.

Cyber threat intelligence (CTI) analysts routinely convert noisy, unstructured security artifacts into standardized, automation-ready representations. Although large language models (LLMs) show promise for this task, existing approaches remain brittle when producing structured CTI outputs and have largely relied on supervised fine-tuning (SFT). In contrast, CTI standards and community-maintained resources define canonical identifiers and schemas that enable deterministic verification of model outputs. We leverage this structure to study reinforcement learning with verifiable rewards (RLVR) for CTI tasks. We introduce \textit{Minerva}, a unified dataset and training pipeline spanning multiple CTI subtasks, each paired with task-specific verifiers that score structured outputs and identifier predictions. To address reward sparsity during rollout, we propose \textit{MinervaRL}, a lightweight self-training mechanism that generates additional verified trajectories and distills them back into the model. Averaged across four backbones and 12 CTI benchmarks, MinervaRL improves the mean score by 15.8 percentage points over the corresponding base models and by 4.3 points over GRPO.

\end{abstract}
